\chapter{Desarrollo}
En esta sección se detalla el proceso de implementación de los diferentes módulos que conforman la aplicación ReMind. Asimismo, se exponen las justificaciones arquitectónicas, las decisiones de diseño técnico y las ventajas operativas que ofrecen las tecnologías seleccionadas para agilizar el desarrollo del sistema.

\section{Desarrollo Backend}
El backend de la aplicación se ha estructurado siguiendo un patrón de arquitectura por capas, organizando el código fuente en paquetes independientes con responsabilidades bien definidas: \textit{entities}, \textit{Mapper}, \textit{repositories}, \textit{services}, \textit{controllers}, \textit{security}, \textit{dto} y \textit{utils}.

A continuación, se detalla la función técnica de cada uno de los componentes:
\begin{itemize}
    \item \textbf{Entity:} Clases de dominio que modelan las entidades del sistema y sus relaciones mediante anotaciones JPA. Cada entidad define su tabla correspondiente, las restricciones de integridad (claves primarias, foráneas, columnas no nulas) y las asociaciones con otras entidades (\texttt{@OneToOne}, \texttt{@ManyToOne}, \texttt{@OneToMany}). Esto permite que Hibernate sincronice automáticamente el esquema de la base de datos con el modelo de objetos.
    \item \textbf{Mapper:} Componentes basados en MapStruct que convierten entidades JPA en DTOs y viceversa, abstrayendo la lógica de mapeo en interfaces declarativas.
    \item \textbf{Repository:} Interfaces que extienden \texttt{JpaRepository}, proporcionando operaciones CRUD estándar y métodos de consulta personalizados sin necesidad de implementación explícita.
    \item \textbf{Service:} Capa donde reside la lógica de negocio. Cada servicio implementa una interfaz separada para facilitar el desacoplamiento y las pruebas unitarias. Aquí se ejecutan las validaciones adicionales, las reglas de negocio y la coordinación entre repositorios.
    \item \textbf{Controller:} Puntos de entrada de la API REST. Cada controlador recibe las peticiones HTTP, valida los parámetros de entrada mediante anotaciones \texttt{@Valid} y delega en la capa de servicio.
    \item \textbf{Security:} Configuración de Spring Security que define la cadena de filtros, las reglas de autorización por ruta y el componente de autenticación JWT.
    \item \textbf{DTO (Data Transfer Object):} Objetos planos que encapsulan los datos que viajan entre cliente y servidor, evitando exponer directamente las entidades JPA y sus relaciones internas.
    \item \textbf{Utils:} Clase de utilidades con funciones auxiliares, principalmente la generación automatizada de nombres de usuario y contraseñas temporales para pacientes y terapeutas.
\end{itemize}

\subsection{Patrones de diseño utilizados}
\begin{itemize}
    \item \textbf{Service Layer:} La lógica de negocio se centraliza en clases de servicio que implementan interfaces específicas. Esto mantiene los controladores ligeros y facilita la prueba unitaria de cada servicio mediante Mockito.
    \item \textbf{DTO:} Cada operación de entrada o salida tiene su propio DTO, de manera que la API expone únicamente los campos necesarios en cada caso y no las entidades completas. Por ejemplo, al crear un paciente se usa \texttt{PacienteInputDTO} y al devolver sus datos se usa \texttt{PacienteOutputDTO}.
    \item \textbf{Repository:} Spring Data JPA permite definir repositorios por herencia de interfaz, con métodos de consulta derivados del nombre del método (como \texttt{findByUsuario()} o \texttt{findByTerapeutaId()}).
    \item \textbf{Lombok:} Las entidades y DTOs se anotan con \texttt{@Data}, \texttt{@AllArgsConstructor} y \texttt{@NoArgsConstructor} para reducir el código repetitivo de getters, setters y constructores.
    \item \textbf{Defensa en profundidad:} La autorización se aplica tanto en la configuración global de Spring Security como en la capa de servicio, donde cada método verifica que el usuario autenticado tenga permiso sobre el recurso solicitado antes de ejecutar cualquier operación.
\end{itemize}

\subsection{Modelo de datos}
El sistema de persistencia se compone de diez entidades JPA que reflejan el modelo relacional de la base de datos. A continuación se describen las más relevantes:

\begin{itemize}
    \item \textbf{Administrador:} Almacena los datos del perfil de administrador. Su tabla contiene los campos básicos de identificación (nombre, apellido, email, teléfono) más las credenciales de acceso (usuario único y contraseña hasheada). El rol se asigna mediante el enum \texttt{TipoUsuario.ADMINISTRADOR}.
    \item \textbf{Terapeuta:} Similar a Administrador pero con un campo adicional de especialidad. Se relaciona con pacientes a través de la tabla intermedia \texttt{PacienteTerapeuta} y con agendas mediante \texttt{AgendaTerapeuta}.
    \item \textbf{Paciente:} Incluye datos personales, la enfermedad diagnosticada, la edad y el nombre del responsable o cuidador. Cada paciente está asociado a un terapeuta mediante una relación \texttt{@ManyToOne} directa en la columna \texttt{terapeuta\_id}.
    \item \textbf{Agenda:} Representa un plan de actividades con un nombre descriptivo. Se vincula a un paciente mediante \texttt{PacienteAgenda} y a un terapeuta mediante \texttt{AgendaTerapeuta}.
    \item \textbf{Juego:} Entidad simple con los campos \texttt{nombre} y \texttt{codigo}, que identifican de forma única cada minijuego disponible en el sistema.
    \item \textbf{JuegoAgenda:} Tabla intermedia que asigna un juego a una agenda con un nivel de dificultad (\texttt{dificultad}, valor entero), una fecha de asignación (\texttt{fechaAsignacion}) y un campo \texttt{realizado} que indica si el paciente ha completado la actividad. Además, incluye \texttt{fechaRealizacion}, que almacena la fecha en que el paciente completó el juego (nulo hasta que se realiza).
    \item \textbf{PacienteJuego:} Registra el historial de juegos realizados por cada paciente, almacenando la fecha de realización y el estado de completado.
\end{itemize}

Todas las tablas se sincronizan automáticamente con la base de datos MySQL mediante la propiedad \texttt{spring.jpa.hibernate.ddl-auto=update}, que crea o actualiza el esquema al iniciar la aplicación.

\subsection{Entorno de desarrollo con Docker}
El entorno de desarrollo del backend se apoya en Docker para la gestión de la base de datos. En lugar de instalar MySQL directamente en el sistema operativo, se define un contenedor Docker a partir de la imagen oficial de MySQL 8.0, configurado mediante el archivo \texttt{docker-compose.yml} situado en la raíz del proyecto. Este archivo especifica las variables de entorno necesarias (usuario \texttt{root}, contraseña \texttt{root1234}, base de datos \texttt{remind}) y monta un volumen (\texttt{mysql\_data}) para garantizar la persistencia de los datos entre reinicios del contenedor. La base de datos se expone en el puerto \texttt{3306} del equipo anfitrión.

El flujo de trabajo diario comienza ejecutando \texttt{docker compose up -d} en la raíz del proyecto para levantar el contenedor MySQL en segundo plano. Una vez que la base de datos está operativa, se inicia el backend con Maven (\texttt{./mvnw spring-boot:run}) y el frontend con Angular CLI (\texttt{ng serve}). Este proceso se ha automatizado mediante el script \texttt{start-all.sh}, que orquesta los tres componentes secuencialmente, permitiendo poner en marcha el entorno completo de desarrollo con una sola orden.

\subsection{DTOs y mapeo}
La comunicación entre la API y la capa de servicio se realiza a través de objetos DTO específicos para cada operación. Se definieron los siguientes grupos:

\begin{itemize}
    \item \textbf{DTOs de entrada:} \texttt{PacienteInputDTO}, \texttt{TerapeutaInputDTO}, \texttt{AgendaInputDTO}, \texttt{AdministradorInputDTO}, \texttt{AssignJuegoDTO}, \texttt{PacienteTerapeutaInputDTO} y \texttt{JuegoInputDTO}. Contienen anotaciones \texttt{@NotBlank}, \texttt{@NotNull} y \texttt{@Size} para validar los datos recibidos.
    \item \textbf{DTOs de salida:} \texttt{PacienteOutputDTO}, \texttt{TerapeutaOutputDTO}, \texttt{AgendaOutputDTO}, \texttt{AdministradorOutputDTO}, \texttt{JuegoAsignadoDTO} y \texttt{JuegoOutputDTO}. Exponen únicamente los campos necesarios para cada vista.
    \item \textbf{DTOs de creación:} \texttt{PacienteCreatedDTO} y \texttt{TerapeutaCreatedDTO}. Se devuelven al registrar un nuevo usuario e incluyen el nombre de usuario y la contraseña generada automáticamente para que el administrador o terapeuta pueda comunicársela al interesado.
    \item \textbf{DTOs de seguimiento:} \texttt{PacienteSeguimientoDTO} y \texttt{TerapeutaSeguimientoDTO}. Contienen datos agregados como el número de juegos completados, el total de juegos asignados y la lista detallada de juegos (con nombre, código, dificultad y estado), usados en el panel de monitorización del terapeuta.
    \item \textbf{DTOs auxiliares:} \texttt{PasswordResetDTO} (devuelve la nueva contraseña tras un reseteo).
\end{itemize}

El mapeo entre entidades y DTOs se realiza mediante MapStruct. Para cada entidad principal existe una interfaz mapper anotada con \texttt{@Mapper(componentModel = "spring")} que define los métodos de conversión. Por ejemplo, \texttt{PacienteMapper}, \texttt{TerapeutaMapper}, \texttt{AgendaMapper}, \texttt{AdministradorMapper} y \texttt{PacienteTerapeutaMapper} convierten sus respectivos DTOs de entrada a entidades y entidades a DTOs de salida, resolviendo automáticamente las relaciones entre tablas mediante \texttt{@Mapping}.

\subsection{Seguridad}
La infraestructura de seguridad se fundamenta en Spring Security y sigue un modelo de autenticación stateless basado en JWT.

\subsubsection{Autenticación JWT}
Cuando un usuario envía sus credenciales a \texttt{POST /api/auth/login}, el \texttt{AuthController} busca secuencialmente en los repositorios de administrador, terapeuta y paciente. Si encuentra una coincidencia, verifica la contraseña con \texttt{PasswordEncoder.matches()} contra el hash BCrypt almacenado. Si la validación es correcta, \texttt{JwtUtil} genera un token HS256 con el nombre de usuario como \textit{subject}, el rol del usuario como \textit{claim} y una expiración de 24 horas.

Para las peticiones posteriores, el cliente debe incluir el token en la cabecera \texttt{Authorization: Bearer \{token\}}. El \texttt{JwtAuthenticationFilter} (un \texttt{OncePerRequestFilter}) intercepta cada petición entrante, extrae el token, lo valida y, antes de establecer la autenticación, verifica que el usuario todavía exista en la base de datos consultando el repositorio correspondiente según su rol. Esto impide que tokens de usuarios eliminados sigan siendo válidos.

\subsubsection{Control de acceso por roles}
Spring Security se configura mediante \texttt{SecurityConfig} con una cadena de reglas que evalúan cada petición HTTP en orden. Las rutas públicas (\texttt{/api/auth/**}) no requieren autenticación. El resto de endpoints están protegidos por rol:

\begin{itemize}
    \item \texttt{/api/administrador/**}: exclusivo para ROLE\_ADMINISTRADOR.
    \item \texttt{/api/terapeuta} (GET): accesible para ADMINISTRADOR y TERAPEUTA.
    \item \texttt{POST, PUT, DELETE /api/terapeuta/**}: solo ADMINISTRADOR.
    \item \texttt{/api/paciente/**} y \texttt{/api/juegos/**}: exclusivo para TERAPEUTA.
    \item \texttt{GET /api/agenda/paciente/**} y \texttt{PUT .../completar}: accesible para PACIENTE y TERAPEUTA.
    \item Cualquier otra ruta requiere autenticación (cualquier rol autenticado).
\end{itemize}

Además, la capa de servicio incorpora sus propias comprobaciones de pertenencia. Por ejemplo, un terapeuta solo puede ver, modificar o eliminar pacientes que tenga asignados; un paciente solo puede consultar su propia agenda. Esta doble verificación (a nivel de ruta y a nivel de servicio) actúa como defensa en profundidad.

\subsubsection{Configuración CORS}
El frontend Angular se ejecuta en \texttt{http://localhost:4200}, por lo que la configuración global de CORS en Spring Boot (mediante \texttt{WebMvcConfigurer}) restringe el origen permitido a esa dirección. Aunque cada controlador incluye \texttt{@CrossOrigin(origins = "*")}, la configuración global centraliza los orígenes permitidos.

\subsubsection{Protección de contraseñas}
Todas las contraseñas se almacenan usando BCrypt mediante la interfaz \texttt{PasswordEncoder} de Spring Security. En ningún momento se guarda la contraseña en texto plano. Cuando se genera una contraseña automática (tras crear un usuario o resetearla), se devuelve en texto plano en la respuesta de la API para que pueda ser comunicada al usuario, pero nunca persiste sin hashear.

\subsection{Servicio de utilidades}
La clase \texttt{Utils} centraliza dos funciones auxiliares empleadas en la creación de usuarios:

\begin{itemize}
    \item \textbf{Generación de nombre de usuario:} Combina el nombre y el apellido del usuario en minúscula, elimina tildes y añade cuatro dígitos aleatorios (por ejemplo, \texttt{juan.perez1234}).
    \item \textbf{Generación de contraseña:} Construye una contraseña temporal combinando las dos primeras letras del nombre en mayúscula con las dos últimas letras del apellido en minúscula, seguidas de un número aleatorio entre 1 y 100 y un carácter especial (por ejemplo, \texttt{NOap42!}).
\end{itemize}

\subsection{Lógica de servicio destacada}
Más allá de las operaciones CRUD estándar, varios servicios implementan comportamiento específico que merece mención:

\subsubsection{Creación de paciente}
Cuando un terapeuta da de alta a un paciente mediante \texttt{POST /api/paciente}, el sistema realiza automáticamente las siguientes operaciones:
\begin{enumerate}
    \item Genera un nombre de usuario y una contraseña temporal mediante \texttt{Utils}.
    \item Codifica la contraseña con BCrypt y crea la entidad \texttt{Paciente} con el terapeuta asignado.
    \item Crea una nueva \texttt{Agenda} con el nombre ``Agenda de \{nombre del paciente\}''.
    \item Crea un registro en \texttt{PacienteAgenda} vinculando el paciente con la agenda.
    \item Crea un registro en \texttt{AgendaTerapeuta} vinculando la agenda con el terapeuta.
    \item Devuelve \texttt{PacienteCreatedDTO} con el usuario y la contraseña en texto plano.
\end{enumerate}
De esta forma, el paciente tiene una agenda lista desde el primer momento y el terapeuta no necesita crearla manualmente.

\subsubsection{Eliminación en cascada}
Al eliminar un terapeuta (\texttt{DELETE /api/terapeuta/\{id\}}), el servicio recorre todos los pacientes asociados y los elimina uno a uno, lo que a su vez dispara la eliminación de sus agendas, relaciones y juegos asignados. Este proceso garantiza que no queden datos huérfanos en la base de datos.

\subsubsection{Verificación de propiedad}
En los servicios de Paciente, Terapeuta y Agenda, cada método que opera sobre un recurso concreto comprueba primero que el usuario autenticado sea el propietario del recurso (o un administrador). Por ejemplo, un terapeuta no puede modificar la agenda de un paciente que no tiene asignado, y un paciente no puede consultar los juegos de otro paciente. Esta verificación se implementa extrayendo el ID del usuario autenticado del \texttt{SecurityContextHolder} y comparándolo con el ID del recurso.

\subsubsection{Seguimiento clínico}
El endpoint \texttt{GET /api/terapeuta/\{id\}/seguimiento} devuelve un \texttt{TerapeutaSeguimientoDTO} que contiene, para cada paciente del terapeuta, el nombre del paciente, el número total de juegos asignados y el número de juegos completados. Esta información se obtiene recorriendo las agendas de cada paciente y contando los registros en \texttt{JuegoAgenda} con el campo \texttt{realizado} a verdadero.

\subsection{API Rest}
A continuación se listan los endpoints implementados, con el método HTTP, la ruta, el rol requerido y una descripción del comportamiento:

\begin{itemize}
\item \textbf{Autenticación}
\begin{itemize}
\item \verb|POST /api/auth/login| (público) Autentica al usuario contra los repositorios de administrador, terapeuta y paciente. Devuelve un token JWT, el rol del usuario, su nombre y su identificador. Si las credenciales no coinciden, devuelve 401.
\end{itemize}

\item \textbf{Administrador}
\begin{itemize}
\item \verb|POST /api/administrador| (ADMINISTRADOR) Registra un nuevo administrador. La contraseña se codifica con BCrypt antes de persistir.
\end{itemize}

\item \textbf{Terapeuta}
\begin{itemize}
\item \verb|GET /api/terapeuta?page=0\&size=10| (ADMINISTRADOR, TERAPEUTA) Lista todos los terapeutas con paginación opcional.
\item \verb|GET /api/terapeuta/{id}| (ADMINISTRADOR, TERAPEUTA) Obtiene un terapeuta por su ID. Los terapeutas solo pueden ver su propio perfil.
\item \verb|POST /api/terapeuta| (ADMINISTRADOR) Crea un terapeuta. Genera usuario y contraseña automáticamente.
\item \verb|PUT /api/terapeuta/{id}| (ADMINISTRADOR) Actualiza los datos de un terapeuta.
\item \verb|DELETE /api/terapeuta/{id}| (ADMINISTRADOR) Elimina un terapeuta y todos sus pacientes asociados en cascada.
\item \verb|GET /api/terapeuta/{id}/seguimiento| (ADMINISTRADOR, TERAPEUTA) Datos de seguimiento clínico: por cada paciente, muestra juegos completados y totales.
\item \verb|GET /api/terapeuta/{id}/agendas| (ADMINISTRADOR, TERAPEUTA) Obtiene todas las agendas gestionadas por el terapeuta.
\item \verb|POST /api/terapeuta/{id}/reset-password| (ADMINISTRADOR, TERAPEUTA) Genera una nueva contraseña, la codifica y la devuelve en texto plano.
\end{itemize}

\item \textbf{Paciente}
\begin{itemize}
\item \verb|GET /api/paciente?page=0\&size=10| (TERAPEUTA) Lista pacientes. Los administradores ven todos; los terapeutas solo los suyos.
\item \verb|GET /api/paciente/{id}| (TERAPEUTA) Obtiene un paciente. El terapeuta solo puede ver sus propios pacientes.
\item \verb|POST /api/paciente| (TERAPEUTA) Crea un paciente. Genera usuario y contraseña, crea una agenda automática y las relaciones correspondientes.
\item \verb|PUT /api/paciente/{id}| (TERAPEUTA) Actualiza los datos de un paciente. Permite cambiar el terapeuta asignado.
\item \verb|DELETE /api/paciente/{id}| (TERAPEUTA) Elimina un paciente y limpia sus relaciones (agenda, juegos asignados).
\item \verb|POST /api/paciente/{id}/reset-password| (TERAPEUTA) Genera una nueva contraseña, la codifica y la devuelve en texto plano.
\end{itemize}

\item \textbf{Agenda}
\begin{itemize}
\item \verb|GET /api/agenda?page=0\&size=10| (TERAPEUTA) Obtiene todas las agendas con paginación.
\item \verb|GET /api/agenda/{id}| (TERAPEUTA) Obtiene una agenda por su ID.
\item \verb|POST /api/agenda| (TERAPEUTA) Crea una agenda y la asocia a un paciente y un terapeuta.
\item \verb|PUT /api/agenda/{id}| (TERAPEUTA) Actualiza el nombre de una agenda.
\item \verb|DELETE /api/agenda/{id}| (TERAPEUTA) Elimina una agenda y sus asociaciones.
\item \verb|GET /api/agenda/paciente/{pacienteId}| (PACIENTE, TERAPEUTA) Obtiene la agenda y los juegos asignados a un paciente. Los pacientes solo ven su propia agenda.
\item \verb|POST /api/agenda/{agendaId}/juego| (TERAPEUTA) Asigna un juego a una agenda con un nivel de dificultad (1--3).
\item \verb|GET /api/agenda/{agendaId}/juegos| (TERAPEUTA) Lista los juegos asignados a una agenda.
\item \verb|PUT /api/agenda/{agendaId}/juego/{juegoId}/completar| (PACIENTE, TERAPEUTA) Marca un juego como realizado y registra la fecha de finalización.
\item \verb|DELETE /api/agenda/{agendaId}/juego/{juegoId}| (TERAPEUTA) Elimina un juego de una agenda.
\end{itemize}

\item \textbf{Juego}
\begin{itemize}
\item \verb|GET /api/juegos| (TERAPEUTA) Obtiene el catálogo completo de juegos disponibles.
\item \verb|GET /api/juegos/{id}| (TERAPEUTA) Obtiene un juego por su identificador.
\item \verb|POST /api/juegos| (TERAPEUTA) Añade un nuevo juego al catálogo.
\item \verb|DELETE /api/juegos/{id}| (TERAPEUTA) Elimina un juego del catálogo.
\end{itemize}
\end{itemize}


\section{Desarrollo Frontend}
La capa de presentación e interfaz interactiva con el usuario se ha desarrollado en Angular. Este framework facilita la estructuración de la interfaz bajo el modelo de Aplicación de una Sola Página (SPA), optimizando la velocidad de respuesta y eliminando las recargas físicas del navegador.

\subsection{Arquitectura del frontend}
La arquitectura del lado del cliente se distribuye según los siguientes principios técnicos de ingeniería de software:
\begin{itemize}
    \item \textbf{Estructura basada en componentes:} Cada módulo visual de la interfaz encapsula su propia lógica de control (TypeScript), plantilla estructural (HTML5) y estilos de diseño independientes (CSS3). Esta modularidad promueve la reutilización de código y minimiza la deuda técnica del proyecto.
    \item \textbf{Organización de directorios:}
    \begin{itemize}
        \item \textbf{components:} Almacena los elementos moleculares y componentes visuales reutilizables de la interfaz, tales como botones, tarjetas y formularios de control.
        \item \textbf{services:} Alberga los servicios de Angular que gestionan la comunicación asíncrona con la API, así como el interceptor HTTP (\texttt{auth.interceptor.ts}) que adjunta automáticamente el token JWT en las cabeceras de cada petición saliente.
        \item \textbf{models:} Define las interfaces y estructuras de datos estáticas que tipifican los objetos compartidos en la aplicación, garantizando la robustez del tipado.
        \item \textbf{pages:} Representa las vistas completas del sistema que se encargan de organizar y orquestar a los diferentes componentes reutilizables.
        \item \textbf{guards:} Contiene los guardianes de ruta encargados de proteger el acceso a las vistas según el estado de autenticación y el rol del usuario (\texttt{terapeuta}, \texttt{paciente} o \texttt{administrador}).
    \end{itemize}
    \item \textbf{Estilos y consistencia visual:} El diseño visual y la adaptabilidad de la interfaz se apoyan en el estándar CSS3 combinado con el framework Bootstrap para agilizar la maquetación geométrica de contenedores, tarjetas, botones y menús. Asimismo, se integra la librería PrimeIcons para unificar el catálogo de iconografía, optimizando la experiencia de usuario y la calidad de la presentación gráfica.
\end{itemize}

\subsection{Comunicación con el Backend}
La transferencia asíncrona de datos entre el cliente Angular y el servidor Spring Boot se realiza mediante el ciclo de peticiones HTTP estructurado bajo los siguientes flujos operativos:

\begin{itemize}
    \item \textbf{Autenticación y ciclo de vida del token:}
    \begin{itemize}
        \item \textbf{Solicitud de Login:} Desde el módulo de acceso, el usuario introduce sus credenciales de seguridad. El servicio asignado captura estos datos y despacha una solicitud asíncrona de tipo POST hacia la API.
        \item \textbf{Validación en el Backend:} El componente \texttt{AuthController} en el servidor intercepta la petición y evalúa la coincidencia de la clave enviada contra el hash criptográfico almacenado en la persistencia.
        \item \textbf{Emisión y almacenamiento de la respuesta:} Tras la validación de las credenciales, el backend transmite el token JWT hacia el cliente, donde el servicio \texttt{auth.service} se encarga de recibirlo y almacenarlo de forma segura en el almacenamiento local del navegador del usuario.
    \end{itemize}
    \item \textbf{Acceso a recursos protegidos:}
    \begin{itemize}
        \item \textbf{Inyección mediante interceptor:} Ante cualquier solicitud de datos protegidos, el interceptor HTTP extrae automáticamente el token almacenado en el navegador y lo adjunta dentro de la cabecera \texttt{Authorization} de la petición saliente.
        \item \textbf{Servicios clientes:} Los servicios especializados inician las llamadas asíncronas orientadas a obtener, crear, modificar o eliminar datos, interactuando de forma limpia con los controladores del backend.
        \item \textbf{Procesamiento y renderizado:} El controlador del backend ejecuta la regla de negocio y transmite una respuesta estructurada en formato JSON que el frontend procesa dinámicamente para actualizar la vista sin necesidad de recargar la página.
    \end{itemize}
\end{itemize}